% Generated from locale/tr/content/first-order-logic/introduction/syntax.tex; do not hand-edit.


\olfileid[tr]{fol}{int}{syn}

\olsection{Sözdizim}

Önce hangi simge dizgelerinin birinci dereceden mantığın
!!{sentence}s ifadeleri sayıldığını kesinleştirmeliyiz. Bunu daha sonra
yapacağız; şimdilik örneklerle ilerleyelim. Birinci dereceden mantığın
temel yapı taşları---söz dağarcığı---iki bölüme ayrılır. İlk bölüm,
belirli şeyler söylemek ya da belirli şeyleri göstermek için
kullandığımız simgelerdir. Şeyleri !!{constant}s aracılığıyla gösterir,
gösterdiğimiz şeyler hakkında !!{predicate}s kullanarak bir şeyler
söyleriz. Örneğin tek bir şeyi göstermek için !!a{constant} olarak
$\Obj a$ kullanabilir ve ardından o şey hakkında
!!{sentence}~$\Atom{\Obj P}{\Obj a}$ ile bir şey söyleyebiliriz.
``$\Obj a$'' ve ``$\Obj P$'' için aklınızda anlamlar varsa
$\Atom{\Obj P}{\Obj a}$ ifadesini doğal dilde bir tümce gibi
okuyabilirsiniz (biçimsel mantığı ilk öğrendiğinizde muhtemelen böyle
yaptınız). Birinci dereceden mantığın bu basit !!{sentence}s ifadeleri
elde edildikten sonra söz dağarcığının ikinci bölümünü, yani mantıksal
simgeleri (bağlaçlar ve niceleyicileri) kullanarak daha karmaşık olanlar
kurabiliriz. Örneğin $(\Atom{\Obj P}{\Obj a}\land
\Atom{\Obj Q}{\Obj b})$ ya da
$\lexists[\Obj x][\Atom{\Obj P}{\Obj x}]$ gibi ifadeler oluşturabiliriz.

Titiz bir metamantıksal inceleme için gereken semantik tanımlarını ve
!!{derivation} sistemlerimizin kurallarını verebilmek üzere, önce neyin
birinci dereceden mantığın !!a{sentence} ifadesi sayıldığını kesin
olarak tanımlamalıyız. Temel fikri anlamak kolaydır: yalnızca
!!{predicate}s ve !!{constant}s kullanarak
$\Atom{\Obj P}{\Obj a}$ gibi bazı basit !!{sentence}s oluşturabiliriz.
Ardından bunlardan bağlaçları ve niceleyicileri kullanarak daha karmaşık
ifadeler oluştururuz. Fakat daha karmaşık !!{sentence}s oluşturmamıza
izin veren kurallar tam olarak nelerdir? Bunlar belirtilmelidir; aksi
halde ``birinci dereceden mantığın !!{sentence} ifadesi''ni yeterince
kesin tanımlamış olmayız. Birkaç sorun vardır. Birincisi, doğru
dizgelerin !!{sentence}s sayılmasını sağlamaktır. İkincisi, bunu bütün
!!{sentence}s hakkında matematiksel ispatlar verebileceğimiz biçimde
yapmaktır. Son olarak ``$!A$ içindeki her $\Obj x$ yerine~$\Obj b$
koy'' gibi, !!{sentence}s üzerinde bazı temel işlemleri de kesin
olarak tanımlamamız gerekecektir. Güçlük, birinci dereceden mantıktaki
niceleyiciler ve !!{variable}s yüzünden bunun nasıl yapılacağının
tamamen açık olmamasıdır. Örneğin
$\lexists[\Obj x][\Atom{\Obj P}{\Obj a}]$ !!a{sentence} sayılmalı
mıdır? Ya $\lexists[\Obj x][\lexists[\Obj x][\Atom{\Obj P}{\Obj x}]]$
ifadesi? ``$(\Atom{\Obj P}{\Obj x}\land
\lexists[\Obj x][\Atom{\Obj P}{\Obj x}])$ içindeki $\Obj x$ yerine
$\Obj b$ koy'' işleminin sonucu ne olmalıdır?

